\documentclass[11pt,a4paper]{article}
\usepackage[utf8]{inputenc}
\usepackage[vietnamese]{babel}
\usepackage{amsmath,amssymb}
\usepackage{geometry}
\geometry{a4paper, total={180mm,260mm}, left=15mm, right=15mm, top=18mm, bottom=18mm}
\usepackage{tikz}
\usetikzlibrary{shapes.geometric, arrows.meta, positioning, calc}
\usepackage{tcolorbox}
\tcbuselibrary{skins,breakable}
\usepackage{xcolor}
\usepackage{multicol}
\usepackage{array}
\usepackage{fontawesome}
\usepackage{fancyhdr}

\pagestyle{fancy}
\fancyhf{}
\renewcommand{\headrulewidth}{0.5pt}
\renewcommand{\footrulewidth}{0.5pt}
\fancyhead[L]{\small\bfseries\color{teal}\faBook\ SỔ TAY KIẾN THỨC TOÁN 3}
\fancyhead[R]{\small\bfseries\color{teal}Lớp Toán Thầy Lĩnh -- loptoanthaylinh.vn}
\fancyfoot[C]{\small\color{gray}-- Sổ tay kiến thức: Bài 68. Tiền Việt Nam --}

\definecolor{primary}{HTML}{005F73}
\definecolor{secondary}{HTML}{0A9396}
\definecolor{accent}{HTML}{EE9B00}
\definecolor{danger}{HTML}{AE2012}
\definecolor{boxbg}{HTML}{F4F6F8}

\begin{document}

% --- BANNER TIÊU ĐỀ ---
\begin{tcolorbox}[
  enhanced,
  colback=white,
  colframe=primary,
  arc=3mm,
  boxrule=1.2pt,
  left=4mm, right=4mm, top=2.5mm, bottom=2.5mm,
  drop shadow=black!10!white
]
  \centering
  {\Large\bfseries\color{primary}\faBookmark\ BÀI 68: TIỀN VIỆT NAM}\\[1.5mm]
  {\small\itshape\color{black!75}Các tờ tiền Việt Nam, nhận biết mệnh giá, kĩ năng đổi tiền và tính tiền thừa khi mua bán trên 1 trang A4}
\end{tcolorbox}

\vspace{-2mm}


\noindent
\begin{minipage}[t]{0.49\textwidth}\vspace{0pt}
\begin{tcolorbox}[
  enhanced,
  colback=blue!4!white,
  colframe=blue!75!black,
  arc=2mm, boxrule=0.8pt,
  left=2mm, right=2mm, top=1.8mm, bottom=1.8mm,
  title={\footnotesize\bfseries\color{white}\faKey\ 1. Các mệnh giá Tiền Việt Nam}
]
{\footnotesize
\begin{itemize}\setlength{\itemsep}{1.5pt}\setlength{\leftskip}{-2.5mm}
  \item Đơn vị tiền tệ của nước ta là \textbf{đồng} (kí hiệu: đ hoặc đồng).
  \item \textbf{Các mệnh giá tiền giấy / cotton nhỏ:}\\
  $100$ đồng, $200$ đồng, $500$ đồng, $1000$ đồng, $2000$ đồng, $5000$ đồng.
  \item \textbf{Các mệnh giá tiền polymer lớn hơn:}\\
  $10\,000$ đồng, $20\,000$ đồng, $50\,000$ đồng, $100\,000$ đồng.
  \item Mỗi tờ tiền có ghi con số mệnh giá, hình hoa văn đặc trưng và chân dung Bác Hồ.
\end{itemize}
}
\end{tcolorbox}
\end{minipage}
\hfill
\begin{minipage}[t]{0.49\textwidth}\vspace{0pt}
\begin{tcolorbox}[
  enhanced,
  colback=teal!4!white,
  colframe=teal!75!black,
  arc=2mm, boxrule=0.8pt,
  left=2mm, right=2mm, top=1.8mm, bottom=1.8mm,
  title={\footnotesize\bfseries\color{white}\faSliders\ 2. Kĩ năng tính tiền và đổi tiền}
]
{\footnotesize
\begin{itemize}\setlength{\itemsep}{1.5pt}\setlength{\leftskip}{-2.5mm}
  \item \textbf{Tính tổng số tiền mua hàng:}\\
  Lấy giá tiền từng món đồ cộng lại với nhau.
  \item \textbf{Tính tiền thừa trả lại:}\\
  $$\text{Tiền thừa} = \text{Tiền khách đưa} - \text{Tiền hàng}$$
  \item \textbf{Đổi tiền tương đương:}\\
  Ví dụ: $1$ tờ $10\,000$ đồng đổi được $2$ tờ $5000$ đồng, hoặc $5$ tờ $2000$ đồng, hoặc $10$ tờ $1000$ đồng.
\end{itemize}
}
\end{tcolorbox}
\end{minipage}

\vspace{-2.5mm}
% --- KHỐI 3: SƠ ĐỒ TƯ DUY TIKZ ---
\noindent
\begin{tcolorbox}[
  enhanced,
  colback=orange!4!white,
  colframe=orange!80!black,
  arc=2mm, boxrule=0.8pt,
  left=2mm, right=2mm, top=1.5mm, bottom=1.5mm,
  title={\footnotesize\bfseries\color{white}\faSitemap\ SƠ ĐỒ TƯ DUY: HỆ THỐNG MỆNH GIÁ TIỀN VÀ BÀI TOÁN MUA BÁN}
]
\centering
\begin{tikzpicture}[
  box/.style={draw, rounded corners=2pt, align=center, font=\scriptsize\bfseries, inner sep=2.5pt, minimum height=0.75cm},
  arr/.style={-{Stealth[length=2mm]}, thick}
]
  \node[box, fill=purple!15, draw=purple!80!black, text width=2.6cm] (m1) at (0, 0) {MỆNH GIÁ NHỎ\\$1000, 2000, 5000$\\Tiền giấy cotton};
  \node[box, fill=teal!15, draw=teal!80!black, text width=2.6cm] (m2) at (3.5, 0) {MỆNH GIÁ LỚN\\$10\,000, 20\,000$\\$50\,000, 100\,000$};
  \node[box, fill=blue!12, draw=blue!80!black, text width=2.6cm] (m3) at (7.0, 0) {ĐỔI TIỀN\\$10\,000 = 2 \times 5000$\\$100\,000 = 2 \times 50\,000$};
  \node[box, fill=orange!15, draw=orange!80!black, text width=2.6cm] (m4) at (10.5, 0) {MUA SẮM\\Tiền hàng = Tổng món\\Tiền thừa = Đưa - Hàng};

  \draw[arr, purple!80!black] (m1) -- (m2);
  \draw[arr, teal!80!black] (m2) -- (m3);
  \draw[arr, blue!80!black] (m3) -- (m4);
\end{tikzpicture}
\end{tcolorbox}

\vspace{-2.5mm}
% --- KHỐI 4: VÍ DỤ & MẸO TRÁNH BẪY ---
\noindent
\begin{minipage}[t]{0.49\textwidth}\vspace{0pt}
\begin{tcolorbox}[
  enhanced,
  colback=green!4!white,
  colframe=green!60!black,
  arc=2mm, boxrule=0.8pt,
  left=2mm, right=2mm, top=1.8mm, bottom=1.8mm,
  title={\footnotesize\bfseries\color{white}\faCheckCircle\ Ví dụ thực hành chuẩn mực}
]
{\footnotesize
\textbf{1. Đếm tiền:}\\
Một ví tiền có: $1$ tờ $50\,000$ đồng, $2$ tờ $20\,000$ đồng và $1$ tờ $5000$ đồng. Tổng số tiền trong ví là:\\
$50\,000 + 20\,000 \times 2 + 5000 = \mathbf{95\,000\text{ đồng}}$.\\
\textbf{2. Bài toán mua sắm:}\\
Bạn Mai mua một quyển truyện giá $25\,000$ đồng và một cái bút giá $8000$ đồng. Mai đưa cô bán hàng tờ $50\,000$ đồng.\\
-- Tiền hàng: $25\,000 + 8000 = 33\,000$ đồng.\\
-- Tiền thừa: $50\,000 - 33\,000 = \mathbf{17\,000\text{ đồng}}$.
}
\end{tcolorbox}
\end{minipage}
\hfill
\begin{minipage}[t]{0.49\textwidth}\vspace{0pt}
\begin{tcolorbox}[
  enhanced,
  colback=red!4!white,
  colframe=red!70!black,
  arc=2mm, boxrule=0.8pt,
  left=2mm, right=2mm, top=1.8mm, bottom=1.8mm,
  title={\footnotesize\bfseries\color{white}\faExclamationTriangle\ Mẹo tránh bẫy trắc nghiệm}
]
{\footnotesize
\begin{itemize}\setlength{\itemsep}{1.5pt}\setlength{\leftskip}{-2.5mm}
  \item \textbf{Bẫy quên đơn vị "đồng":}\\
  Đáp số của bài toán tính tiền bắt buộc phải ghi rõ tên đơn vị là \textbf{"đồng"}, không được bỏ trống!
  \item \textbf{Bẫy nhầm tờ $20\,000$ và $500\,000$:}\\
  Hai tờ polymer này đều có tông màu xanh lơ gần giống nhau, nhưng lớp 3 chỉ học đến mệnh giá cao nhất là $\mathbf{100\,000\text{ đồng}}$.
\end{itemize}
}
\end{tcolorbox}
\end{minipage}


\end{document}
